Dans le plan muni d'un repère orthonormal, on considère les points \mbox{P~$(4;0)$} et \mbox{K~$(1;0)$}.

On considère le réel $x$ et on note M le point de coordonnées \mbox{$(x;3)$.}
\begin{center}
	\begin{tikzpicture}[xscale=1,yscale=1,rotate=0,every node/.style={scale=1,rotate=0}]
	\clip(-0.6,-0.6)rectangle(4.6,3.6);
	\draw[blue!30,thick](0,0)grid(4.6,3.6);
	\draw[line width=0.3mm,->, >=triangle 45](-0.6,0)--(4.6,0);
	\draw[line width=0.3mm,->, >=triangle 45](0,-0.6)--(0,3.6);
	\foreach \i in {1,2,...,4}{\draw(\i,-0.2)node{\textbf{\i}};}
		\foreach \j in {1,2,3}{\draw(-0.2,\j)node{\textbf{\j}};}
		\draw(-0.2,-0.2)node{\textbf{0}};
	\draw[fill=black](1,0)circle(0.04)node[above right]{\textbf{K}}(4,0)circle(0.04)node[above right]{\textbf{P}}(2.3,3)circle(0.04)node[above]{\textbf{M}};
\end{tikzpicture}
\end{center}

\begin{enumerate}
	\item Donner les coordonnées du vecteur \mbox{$\vectt{KP}$} ainsi que sa norme.

	\item Exprimer en fonction de $x$ les coordonnées du vecteur \mbox{$\vectt{KM}$} ainsi que sa norme.

	\item Montrer que le produit scalaire \mbox{$\vectt{KP}\cdot\vectt{KM}$} est égal à \quad \mbox{$3x-3$.}

	\item Montrer que si l'angle $\widehat{\text{PKM}}$ est égal à \mbox{$\dfrac\pi 3$}, alors le réel $x$ est solution de l'équation
	\[(E)\quad \sqrt{(x-1)^2+9}=2x-2\]

	\item Vérifier que le réel \mbox{$1+\sqrt3$} est solution de l'équation (\textit{E}).
\end{enumerate}
